\documentclass[11pt]{article}
\usepackage[a4paper, top=18mm, bottom=18mm, left=20mm, right=20mm]{geometry}
\usepackage[T2A]{fontenc}
\usepackage[utf8]{inputenc}
\usepackage[ukrainian]{babel}
\usepackage{graphicx}
\usepackage[export]{adjustbox}
\usepackage{amsmath}
\usepackage{amssymb}
\graphicspath{{/app/Quiz/Scripts/2023/ProgAll/15BinTree/}{/app/Quiz/Scripts/2021/Mod2021_3/BinTree/}{/app/Quiz/Scripts/2020/Mod2020_4/BinTree/}}
\setlength{\parindent}{0pt}
\setlength{\parskip}{6pt}
\pagestyle{empty}
\begin{document}
%begin
{\large\textbf{Контрольна робота}}

\textbf{Варіант \No\,1}\hfill \textit{ПІБ:}\,\underline{\hspace{60mm}}
\vspace{4mm}

%beginitem
1. %goodpath:Scripts/2018/Mod2018_1/03-good.txt
Відмітити все, що є коректним оператором С++:
( 1 ) bool f(int a);

( 2 ) int f(int a,b);

( 3 ) cout<<3.4;

( 4 ) int n; bool n;
\vspace{5mm}
%enditem

%beginitem
2. Що буде виведено за виконання фрагменту коду?

%code
cout << (12 % 12 % 12 % 12);
%code
\vspace{5mm}
%enditem

%beginitem
3. Що буде виведено за виконання фрагменту коду?

a,b,c=6,2,7
def h(a):
global c    a=3
    b*=2
    c=4
    return a+b+c

a,b,c=9,4,3
print(h(a),a,b,c)
\vspace{5mm}
%enditem

%beginitem
4. Що буде виведено за виконання фрагменту коду?

%code
print(8 > 4 <= 8.0 is True)
%code
\vspace{5mm}
%enditem

%beginitem
5. Що буде виведено за виконання фрагменту коду?
%code
a, b, c = 6, 8, 7
def h(a, b, c):
    print(a, b, c, end=" ")

h(b=0, c=2, 1)
print(a, b, c)
%code
\vspace{5mm}
%enditem

%beginitem
6. %goodpath:Scripts/2019/Mod2019_1/Lex/01-good.txt
Відмітити все, що є однією правильною лексемою:
( 1 ) False

( 2 ) '\'

( 3 ) py

( 4 ) 8,3
\vspace{5mm}
%enditem

%beginitem
7. Що буде виведено за виконання фрагменту коду?

print('R', end=' ')
g=[10,11,12,13,14]
print(g.pop(1), end=' ')
print(g)

\vspace{5mm}
%enditem

%beginitem
8. Що буде виведено за виконання фрагменту коду?

%code
9**0.5*-2
%code
\vspace{5mm}
%enditem

%beginitem
9. 

\begin{center}
	\adjustimage{max size={0.8\linewidth}{0.8\paperheight}}
	{../15BinTree/trees_exp/178.png}
\end{center}
Указати послідовність міток вершин дерева, зображеного на рис., що утвориться в результаті обходу дерева в прямому порядку. Мітки записувати через пробіл.\par Алгоритм починає роботу з кореня (на рис. це верхня вершина).
%answer:178 прямому
\vspace{5mm}
%enditem

%beginitem
10. %goodpath:Scripts/2019/Mod2019_1/Lex/03-good.txt
Відмітити все, що є коректним оператором С++:
( 1 ) 3.14%2

( 2 ) if a<b: a=0 else: b=0

( 3 ) i,j=j,j

( 4 ) x+3=5
\vspace{5mm}
%enditem

%end
\newpage
\end{document}

